\documentclass[11pt,a4paper]{article}
\usepackage{amsmath,amssymb,amsthm}
\usepackage[margin=2.5cm]{geometry}
\usepackage{graphicx}
\usepackage{xcolor}
\usepackage{booktabs}
\usepackage{tabularx}
\usepackage{longtable}
\usepackage[export]{adjustbox}
\usepackage{enumitem}
\usepackage{microtype}
\usepackage{needspace}
\usepackage[section]{placeins}
\usepackage[colorlinks=true,linkcolor=blue!55!black,citecolor=blue!55!black,
            urlcolor=blue!55!black,bookmarks=true,bookmarksnumbered=true,
            unicode=true]{hyperref}

\widowpenalty=10000
\clubpenalty=10000

\theoremstyle{plain}
\newtheorem{theorem}{Theorem}[section]
\newtheorem{lemma}[theorem]{Lemma}
\newtheorem{corollary}[theorem]{Corollary}
\newtheorem{proposition}[theorem]{Proposition}
\newtheorem{conjecture}[theorem]{Conjecture}
\theoremstyle{definition}
\newtheorem{definition}[theorem]{Definition}
\newtheorem{example}[theorem]{Example}
\theoremstyle{remark}
\newtheorem{remark}[theorem]{Remark}

\newcommand{\Z}{\mathbb{Z}}
\newcommand{\R}{\mathbb{R}}
\newcommand{\T}{\mathbb{T}}
\newcommand{\floor}[1]{\left\lfloor #1 \right\rfloor}
\newcommand{\wn}[1]{\lVert #1 \rVert_{N}}
\newcommand{\wnc}[1]{\lVert #1 \rVert}
\newcommand{\lrc}{\mathrm{LRC}}
\newcommand{\taust}{\mathrm{TAU}}
\DeclareMathOperator{\ord}{ord}
\DeclareMathOperator{\cls}{cls}
\DeclareMathOperator{\Vol}{Vol}

\graphicspath{{.}}

\hypersetup{
  pdftitle={Every Natural Strengthening Is False: Type Mismatch in
            Lossless Reformulations of the Lonely Runner Conjecture},
  pdfauthor={Anonymous},
  pdfsubject={Lonely Runner Conjecture; reformulations; covering radius;
              zonotopes; Ehrhart theory; negative results},
  pdfkeywords={Lonely Runner Conjecture, lossless reformulation,
               type mismatch, covering radius, LR-zonotope, Ehrhart
               h*-polynomial, arithmetic matroid, exact witnesses,
               branch and bound certification}
}

\begin{document}

\tableofcontents
\newpage

% ======================================================================
% Section 1 — Abstract and Introduction
% ======================================================================

\begin{abstract}
\noindent
The Lonely Runner Conjecture ($\lrc_n$) admits many \emph{lossless}
reformulations---finite, geometric, dynamical, or arithmetic rewrites that
are exactly equivalent to it. This paper is about a structural fact that
emerged from a multi-year program of attacking the conjecture through such
reformulations: \emph{every natural strengthening of a lossless
reformulation---every version of it that would open a new toolkit---is
false}. A reformulation equivalent to $\lrc$ carries the full difficulty of
$\lrc$; a strictly stronger version is false, and the falsity is visible
at small $n$ with exact rational witnesses. We exhibit this type mismatch
in four independent settings: the finite pair-sum covering reformulation,
whose classification toolkit provably emits only closure, invariance, and
exhaustiveness statements (a full proof-inventory audit) and whose single
margin-shaped instrument, a transfer-matrix model of the solution cascade,
is closed by theorem (its constrained transfer matrices are partial
permutations with spectral radius exactly $0$ or $1$); the
character-rigidity route, whose orbit-invariance laws are falsified
exactly and whose marginal envelopes saturate precisely where the missing
margin would be needed; and the Lonely Runner zonotope, whose covering
radius satisfies a strictly stronger inequality than $\lrc$ that is false
from $n=5$ on, with exact witnesses
$\rho \ge 97/288 > 1/3$ at $v=(1,2,3,4,5)$,
$\rho \ge 29/80 > 5/14$ at $v=(1,2,3,4,5,6)$, and
$\rho \ge 4/11 > 5/14$ at the non-extremal instance $v=(1,2,3,4,5,7)$,
while its Ehrhart $h^*$-polynomials are real-rooted and log-concave
through $n=10$ with strictness decaying like $\Theta(1/n)$ and no
correlation with covering margins.
Alongside the negative result we record the program's one bounded positive
theorem---the sum-only equality $M(V)=\max_{\{u,w\}}\tau(u,w)$, the
rigid-zone translate-class classifications at the primes
$7,11,13,17,19,37$, and the composite-side structure at $p^2$ and $pq$,
stated with their hypothesis ledger---and a finite witness archive with a
certification ladder: exact deepest-hole classification at $n=3$,
independent away-certified and branch-and-bound certificates at $n=4$
(newly including $v=(1,2,3,16)$, $\rho=9/34$), exact refutations at
$n=5,6$, and, from this session's independent branch-and-bound method,
certified deepest-hole equality at $n=5$, $m=20,25,30$
($\rho=\delta$ exactly at each) together with one serious attempt at the
remaining $m=10$ case, whose unresolved residual has volume below
$6\times10^{-7}$ with exact interior depths strictly below the claimed
value.
The paper is a negative-results paper and claims no progress on the
conjecture itself: it proves no new cases of $\lrc$, and the surviving
equivalent forms are as hard as they ever were. What it contributes is the
obstruction's anatomy: why four natural toolkits, each correct and each
productive on its own terms, cannot cross the same gap.
\end{abstract}

\section{Introduction}\label{sec:intro}

\subsection{The conjecture and the reformulation habit}

Let $n\ge 2$ and let $V=(v_1,\dots,v_n)$ be distinct positive integers
with $\gcd(V)=1$. Runners start together at the origin of the unit circle
$\T=\R/\Z$; runner $p$ occupies $v_p t \bmod 1$ at time $t$. Write
$\wnc{x}$ for the distance of a real number to the nearest integer. The
\emph{loneliness} of the stationary observer at time $t$ is
$g(t)=\min_{1\le p\le n}\wnc{v_p t}$, and $M(V)=\max_{t\in\T}g(t)$. The
Lonely Runner Conjecture, posed by Wills in 1967 and independently by
Cusick in the language of view obstruction
\cite{PS-survey,Wills}, asserts $M(V)\ge 1/(n{+}1)$ for every $V$. It is
known for at most seven runners \cite{PS-survey}; the frontier
(eight through fifteen runners) is entirely computer-assisted
\cite{Rosenfeld8,Trakulthongchai,Sungkawichai,Allikvere,MSS}.

Attacking such a conjecture through equivalent rewrites is an old habit.
The reformulation landscape around the Lonely Runner problem is unusually
rich: Wills' original inhomogeneous-approximation form; Cusick's view
obstruction; Rifford's continuous covering problem in dimension $n-2$,
equivalent to a time-bound version of the problem \cite{Rifford}; and, from
the program this paper reports on, a finite form---the pair-sum lattice
reduction---under which the conjecture becomes a family of covering
statements about explicit cyclic groups \cite{PS-survey,prog-artifacts}.
Each rewrite is \emph{lossless}: it holds exactly when $\lrc_n$ holds,
with no residual error term. Each was adopted in the hope that its native
toolkit---covering combinatorics, spectral theory, harmonic analysis,
polyhedral geometry, Ehrhart theory---could move the conjecture.

This paper records, in a form we hope is useful beyond this one problem,
what that hope runs into. The headline is a sharp negative observation,
verified independently four times:

\begin{quote}
\emph{LRC admits many lossless reformulations, and every natural
strengthening of them---every reformulation that would open a new
toolkit---is false. This is not a limitation of one technique; it is a
structural fact, exhibited here in four independent settings with finite
exact witnesses.}
\end{quote}

The mechanism, which we call \emph{type mismatch} and develop in
\S\ref{sec:tm}, has three faces. First, a reformulation equivalent to
$\lrc$ is as hard as $\lrc$; losslessness guarantees that nothing is
discarded, so nothing is easier. Second, the statements that would make a
reformulation tractable---a uniform margin, a spectral gap, a saturated
bound with slack, a monotone scale law---are strictly stronger than $\lrc$
wherever they are false at the extremal instances, and they \emph{are}
false there. Third, the instruments that a lossless reformulation
naturally supports (classifications, invariances, exact counts) provably
compose into statements of the wrong type: they close everything except
the one gap that matters. All four faces were verified by computation,
with exact arithmetic, and the witnesses are small, checkable, and
archived (\S\ref{sec:witness}, Appendix~\ref{sec:prov}).

\subsection{The four instances in one paragraph each}

\emph{Instance 1: the finite covering reformulation}
(\S\ref{sec:tm-covering}). The pair-sum equality theorem reduces
$M(V)$ exactly to a maximum over cyclic groups $\Z_{v_u+v_w}$
(\S\ref{sec:positive}), turning $\lrc_n$ into a covering question about
explicit bad sets. The program's classification toolkit (size laws, moat
counts, class-group reductions, lift laws) closed $75.6\%$ of a
$3{,}712{,}576$-instance corpus outright and localized the rest to one
named class---but a statement-by-statement audit of the full proof
inventory found \emph{zero} proved statements of the missing type (a
uniform growth-in-$k$ margin over a trivial baseline); every proved
statement is a closure, invariance, exhaustiveness, or machine
verification, and the one candidate that reads like an exception (the cut
law's ``$k$-cancellation'') is a $k$-\emph{invariance}---the exact
opposite of a growth-in-$k$ margin.

\emph{Instance 2: the transfer-matrix route} (\S\ref{sec:tm-transfer}).
The single instrument shaped to produce the missing type---a
transfer-matrix or automaton model of the solution cascade, whose natural
output is a spectral quantity uniform in scale---was executed as a proof
attempt under a pre-committed success criterion and closed by theorem:
every constrained transfer matrix of the lossless cascade is a partial
permutation with spectral radius exactly $0$ or $1$; the cyclic product is
diagonal, its trace is the survivor count itself, and at the six
reflection-symmetric committed families the survival fraction is exactly
$1$ forever. Losslessness forbids generated decay: bijections preserve
measure, deletion filters without dissipating, and the reformulation is
bounded by its own fidelity.

\emph{Instance 3: character rigidity} (\S\ref{sec:tm-char}). The
harmonic-analytic route phrased the obstruction in terms of a closed
system of ten characters. Its natural strengthenings---orbit-invariance
of singleton or pairwise marginal profiles---are \emph{false}, falsified
by exhaustive exact computation; the marginal envelope saturates (is
full) exactly at the configurations where the margin is needed, and the
only invariant sparsity is joint. The trivial bound sits exactly at the
cascade threshold: beating it by the factor the obligations require
($\sim\!250$ at every rung) is precisely the open margin, stated twice.

\emph{Instance 4: the Lonely Runner zonotope}
(\S\ref{sec:tm-zono}). The quotient torus $\R^n/(\R v+\Z^n)$ carries a
covering radius $\rho(v)$ in the quotient $\ell_\infty$ norm, and the
Lonely Runner statement for the center class is
$\delta(v)\le (n{-}1)/(2(n{+}1))$ with $\delta\le\rho$ always. The
covering-radius form---\emph{every} hole within the threshold, not just
the center---would open the geometry-of-numbers toolkit. It is false from
$n=5$: at the extremal harmonic instance $v=(1,2,3,4,5)$ there is an
exact rational point at depth $97/288>1/3$, at $v=(1,2,3,4,5,6)$ one at
$29/80>5/14$, and at the non-extremal $v=(1,2,3,4,5,7)$ one at
$4/11>5/14$. Meanwhile the Ehrhart $h^*$-polynomials of the zonotopes
are beautifully behaved---real-rooted, log-concave, Newton, uniformly,
through $n=10$ in both core families---but their strictness decays like
$\Theta(1/n)$ and has no measurable correlation with the covering
margins: the coefficient-level toolkit cannot see the metric hole depth
at all.

\subsection{What this paper is and is not}

The paper is deliberately honest about its shape. It is a \emph{sharp
negative result}---a whole class of approaches to a famous problem, and
the exact reason each fails---plus \emph{one bounded positive theorem}
carried over from the companion monograph (the equality theorem, the
rigid-zone classifications, and the $p^2$/$pq$ composite structure,
\S\ref{sec:positive}) and a \emph{finite witness archive} with an explicit
certification ladder (\S\ref{sec:witness}). Negative results of this kind
are genuinely valuable when they come with exact computational evidence
and a mechanism; the mechanism here (losslessness versus strengthening)
applies verbatim to any problem with a comparable reformulation culture.

The boundary of the claims is as follows. The paper does not solve,
reduce, or weaken the Lonely Runner Conjecture; no new cases of $\lrc$
are proved at any $n$; the rung results of
\cite{Rosenfeld8,Trakulthongchai,Sungkawichai,Allikvere} are untouched.
The equivalent forms that survive our audits (the $\delta$-form of the
covering statement, the pair-sum formulation) are exactly as hard as
$\lrc$, and we claim no leverage on them. The falsified strengthenings
are stronger than $\lrc$ by construction, so their falsity implies
nothing about $\lrc$ itself---which remains open, tight at the harmonic
instances, and apparently indifferent to every toolkit this program could
bring to bear. What the archive adds is certainty about where \emph{not}
to dig: four doors, each closed by exact computation or by theorem, each
closing for the same reason.

\subsection{Roadmap and conventions}

\S\ref{sec:prelim} fixes notation and the four equivalent forms of the
covering statement, with the proofs of equivalence deferred to the
companion monograph where they are long. \S\ref{sec:tm} states the
type-mismatch mechanism and develops the four instances.
\S\ref{sec:positive} summarizes the positive core with exact statements.
\S\ref{sec:witness} presents the witness archive and the certification
ladder, including the new certificates of this session.
\S\ref{sec:concl} concludes. Appendix~\ref{sec:hstar} records the
$h^*$-observations (the extension to $n=9,10$ is new here).
Appendix~\ref{sec:prov} is the provenance ledger: every number cited in
this paper is re-asserted by independent exact computation in a
persisted, machine-checkable run (225 items, zero failures), and the
certification scripts, budgets, and outputs are listed.

All computations are exact (rational arithmetic) unless explicitly marked
as a float enclosure; certification levels are labeled \emph{exact}
(complete exact proof), \emph{certified} (complete proof with an
independent numerical outer loop), or \emph{consistent} (search-level
evidence only), following the ladder fixed in the audit of the underlying
program records \cite{prog-artifacts}.

% ======================================================================
% Section 2 — Preliminaries
% ======================================================================

\section{Preliminaries}\label{sec:prelim}

\subsection{Notation}

Throughout, $n\ge 3$ is the number of speeds, $V=(v_1,\dots,v_n)$ has
distinct positive integer entries with $\gcd(V)=1$ and $v_1=1$ when a
coordinate model is in use, and $\wnc{x}=\min(x-\floor{x},
\floor{x}+1-x)$ is the distance to the nearest integer. We write
$\lambda(V)=\max_{t\in\T}\min_p \wnc{v_p t}$ for the optimal loneliness
and $M(V)$ for the same quantity; the two notations are interchangeable
and we keep $\lambda$ in metric contexts and $M$ in the
pair-sum context. The threshold is
$\mathrm{thr}_n=(n-1)/(2(n+1))=1/2-1/(n+1)$. Two speed families recur:
the \emph{harmonic} family $V=(1,2,\dots,n)$ and the \emph{rung-2}
family $V=(1,\dots,n-1,2n)$. For these, $\lambda$ is classical:
$\lambda(1,\dots,n)=1/(n+1)$ exactly (the harmonic family is extremal,
with zero margin against the conjecture), and
$\lambda(1,\dots,n-1,2n)=2/(2n+1)$ exactly (the rung-2 family is not
extremal; its margin is $1/((n+1)(2n+1))$). Both identities are
re-asserted exactly for $n\le 10$ in the provenance run of
Appendix~\ref{sec:prov}.

\subsection{Four forms of the covering statement}

Write $\delta(V)=\tfrac12-\lambda(V)$ and let $\mathbf{1}=(1,\dots,1)$.
The following forms of $\lrc_n$ are used interchangeably in the
literature and in this program, and it is essential to keep their
strengths apart.

\begin{enumerate}[leftmargin=1.6em,itemsep=2pt]
\item \textbf{Loneliness form.} $\lambda(V)\ge 1/(n+1)$ for all $V$.
\item \textbf{Center-depth form ($\delta$-form).}
$\delta(V)\le \mathrm{thr}_n$, where
$\delta(V)=\operatorname{dist}_\infty\!\bigl(\R v,\;\Z^n+\tfrac12\mathbf{1}\bigr)$
is the $\ell_\infty$ distance from the trajectory line to the shifted
lattice. The identity
$\delta=\tfrac12-\lambda$ is elementary: for every $t$,
$\max_j \wnc{\tfrac12-v_jt}=\tfrac12-\min_j\wnc{v_jt}$.
\item \textbf{Quotient-torus form.} Let $V=\R^n/\R v$, let
$\pi\colon\R^n\to V$ be the projection, $\Lambda=\pi(\Z^n)\cong\Z^{n-1}$,
and let $\|\cdot\|_V$ be the quotient norm
$\|c\|_V=\operatorname{dist}_\infty(x,\R v)$ for any lift $x$ of $c$.
Then $\delta(V)=\operatorname{dist}_V(c^*,\Lambda)$ where
$c^*=\pi(\tfrac12\mathbf{1})$ is the \emph{center class}, a $2$-torsion
point whose coordinates are $\tfrac12$ on even speeds and $0$ on odd
speeds. The $\delta$-form reads: the center class has quotient-norm
depth at most $\mathrm{thr}_n$ in every $V$.
\item \textbf{Covering-radius form ($\rho$-form).}
Let $\rho(V)=\max_{c\in V}\operatorname{dist}_V(c,\Lambda)$ be the
covering radius of $\Lambda$ in the quotient norm. The $\rho$-form of the
conjecture asserts $\rho(V)\le\mathrm{thr}_n$ for all $V$.
\end{enumerate}

Forms (1)--(3) are equivalent to $\lrc_n$: this is the content of the
shift identity above together with the standard quotient-torus dictionary,
and the equivalence is asserted computationally, exactly, on every
instance of the program (the geometric pipeline is validated end to end
in the runs of Appendix~\ref{sec:prov}). Form (4) is \emph{strictly
stronger}. Since the maximum defining $\rho$ ranges over all classes and
$\delta$ is the depth of one class, $\delta(V)\le\rho(V)$ always; hence
$\rho\le\mathrm{thr}_n$ implies the $\delta$-form and thus $\lrc_n$. The
converse fails, and the failure is not an artifact of small examples: it
begins exactly at $n=5$ and persists (\S\ref{sec:tm-zono},
\S\ref{sec:witness}). At the harmonic instances the $\delta$-form is
tight ($\delta=\mathrm{thr}_n$), so the $\rho$-form has no slack to
absorb the difference between the center and the deepest hole: the
two statements coincide on all tested instances for $n\le 4$---which is
why the identification was easy to believe---and separate from $n=5$ on,
with exact witnesses.

For computations we use two further standard facts, both proved and
machine-validated in the program's records \cite{prog-artifacts}. First,
in \emph{$c$-coordinates} $c_j=x_j-v_jx_1$ ($j=2,\dots,n$; valid when
$v_1=1$) the lattice is exactly $\Z^{n-1}$ and the quotient norm is the
maximum of the finitely many linear forms
\[
\phi^{1j}=\frac{c_j}{1+v_j},\qquad
\phi^{ij}=\frac{v_jc_i-v_ic_j}{v_i+v_j}\quad(2\le i<j),
\]
each of $\ell_1$-norm at most $1$; hence
$\|c\|_V\le\|c\|_\infty$ and Lipschitz certification needs no
norm-equivalence constant. Second, the depth function admits the
one-parameter \emph{runner formulation}
\begin{equation}\label{eq:runnerform}
D(c)\;=\;\min_{\sigma\in[0,1)}\;\max\Bigl(\wnc{\sigma},\;
\max_{j\ge 2}\wnc{c_j-v_j\sigma}\Bigr),
\end{equation}
which is the distance of $c$ to $\Lambda$ on the quotient torus; this is
the form used by the branch-and-bound certifier of \S\ref{sec:witness},
and it was validated against the independent linear-form/lattice-search
implementation on random rational points for every instance used in this
paper (forty points per instance, exact agreement throughout).

\subsection{The Lonely Runner zonotope and its Ehrhart theory}

The \emph{LR-zonotope} of $V$ is $Z(V)=\pi([0,1]^n)\subset V$; it is an
$(n{-}1)$-dimensional lattice zonotope with respect to $\Lambda$, with
normalized volume $\sum_j v_j$. Its lattice-point geometry encodes the
covering problem above: holes of the $\Lambda$-covering correspond to
classes of large quotient depth. Its Ehrhart theory is arithmetic and
clean \cite{BeckRobins,dAMoci,StanleyEC}. Writing $d=n-1$, the Ehrhart
polynomial of $Z(V)$ is
\begin{equation}\label{eq:ehr}
\mathrm{Ehr}_{Z(V)}(t)\;=\;\sum_{k=0}^{d} e_k\,t^k,\qquad
e_k\;=\;\sum_{\substack{S\subseteq[n]\\|S|=k}}\gcd\!\bigl(v_{S^c}\bigr),
\qquad e_d=\sum_j v_j,
\end{equation}
with $\gcd$ of the empty set $0$, and the $h^*$-polynomial is obtained
from the Eulerian expansion
$\sum_{t\ge0}\mathrm{Ehr}(t)z^t=h^*(z)/(1-z)^{d+1}$; explicitly
$h^*(z)=\sum_k e_k\,(1-z)^{d-k}A_k(z)$ with $A_k$ the Eulerian
polynomials. Formula \eqref{eq:ehr} was derived in the program's records
and verified in two independent ways on every instance used here (a DFS
coset count and the standard Ehrhart inversion;
Appendix~\ref{sec:prov}). The gcd structure is that of the arithmetic
matroid of the projected basis \cite{dAMoci}, which is why the
arithmetic-matroid literature is the natural home for the residual
positive question in Appendix~\ref{sec:hstar}.

\subsection{The certification ladder}

Every classification claim in this paper carries one of three labels.
\emph{Exact}: a complete proof in rational arithmetic (the $d=2$
arrangement enumerations; all refuting witnesses, which are single
evaluations of $D$ at rational points). \emph{Certified}: a complete
proof in which an outer numerical loop (a grid or a branch-and-bound
frontier) is combined with an exact inner analysis and a rigorous
one-sided error bound (the $d=3$ away-certifications; the
branch-and-bound certificates of \S\ref{sec:witness}). \emph{Consistent}:
search-level evidence only (grid maxima equal to the claimed value,
multi-start optimization finding nothing above it), with no completeness
claim. The labels were introduced by the audit of the program's earlier
records precisely because enclosure-level evidence had been presented as
proof; the ladder is applied uniformly here, including to the results of
this session.

% ======================================================================
% Section 3 — The type-mismatch mechanism and four instances
% ======================================================================

\section{The type-mismatch mechanism and four instances}\label{sec:tm}

\subsection{The mechanism}\label{sec:tm-mech}

Call a statement $R$ a \emph{lossless reformulation} of $\lrc_n$ if $R$
holds for all speed vectors at rung $n$ if and only if $\lrc_n$ does, with
no slack in either direction. Call a statement $S$ a \emph{toolkit
strengthening} of $R$ if (i) $S$ implies $R$, strictly, and (ii) the
strictness is of the kind that a named toolkit can exploit: a uniform
margin over a trivial baseline, a spectral gap, an interior rate, a
monotone law in the scale parameter, a saturated bound with slack. The
distinction is not aesthetic; it is the difference between statements
that classification, spectral, harmonic, or polyhedral instruments can
consume and statements they cannot.

Three observations, jointly verified four times in this program, make up
the mechanism.

\begin{enumerate}[leftmargin=1.6em,itemsep=2pt]
\item \textbf{Losslessness is as-hard-ness.} A lossless reformulation
carries the full difficulty of $\lrc_n$ by construction. No instrument
that respects the equivalence---that is, no instrument whose outputs are
invariant under the lossless rewriting---can distinguish the
reformulation from the original.
\item \textbf{The strengthening is false where it matters.} The
toolkit-relevant strengthenings are strictly stronger than $\lrc$, and at
the extremal instances (the harmonic and rung-2 families, the
reflection-symmetric cascade families, the tight covering
configurations) the extra margin they demand does not exist. When the
demand is a definite statement, it is false, and small $n$ exhibits the
falsity with exact witnesses.
\item \textbf{The instruments land on the closed side.} The instruments a
lossless reformulation naturally supports---exact counts,
classifications, invariances, exhaustiveness checks---provably compose
into statements that close everything \emph{except} the missing margin.
In the setting of Instance 1 this is a theorem-level audit finding; in
Instance 2 it is a theorem; in Instances 3 and 4 it is measured, with the
saturation and decay records below.
\end{enumerate}

The four subsections that follow develop one instance each. They are
independent: different reformulations, different hoped-for toolkits,
different witnesses. What they share is the mechanism, and we state it
once, here, as the paper's organizing claim: \emph{in each setting, the
falsity of the strengthening is not an accident of the example but a
consequence of the losslessness that made the reformulation faithful.}

\subsection{Instance 1: the finite covering reformulation}\label{sec:tm-covering}

The pair-sum equality theorem (\S\ref{sec:positive}) reduces
$M(V)$ to a finite maximum over cyclic groups: $M(V)=\max_{\{u,w\}}
\tau(u,w)$, where $\tau$ is the optimal loneliness on the pair-sum
lattice of denominator $N=v_u+v_w$. Under this reduction $\lrc_n$
becomes the statement that for every $V$ some pair's bad sets
$B_w=\{k\in\Z_N:(n{+}1)\wn{wk}<N\}$ fail to cover $\Z_N$: a finite
covering problem, exact, and apparently made for the covering toolkit.

The toolkit delivered. The program proved an exact size law for the bad
sets, a polygonal witness lemma, moat-counting and coincidence-counting
lemmas, a lift law, and a class-group reduction at prime moduli; it
classified the unit covering configurations completely at the rigid
primes (Theorem~\ref{thm:rigid}); and on the full corpus of
$3{,}712{,}576$ primitive $5$-sets with speeds $\le 56$ it closed
$75.6\%$ of instances by proved lemmas alone ($78.2\%$ with the
coincidence lemma) and all of them by per-modulus verification,
localizing the residual obstruction to a single named class
\cite{prog-artifacts}. What it did not deliver, at any point in the
program's multi-year run, is a single statement of the type the residual
obstruction requires: a \emph{uniform growth-in-$k$ margin over a trivial
or measured baseline}.

This is a verifiable claim, and it was verified by audit: the full proof
inventory of the program---$34$ named results across the pair-sum,
group-ring, tensor, and cascade sections---was classified
statement-by-statement. Every proved statement is a closure, an
invariance, an exhaustiveness, or a machine verification. The one
candidate that reads like an exception, the cut law's
``$k$-cancellation,'' resolves as a $k$-scale \emph{invariance} of the
composed bound: it cancels exactly the $k$-dependence that a margin
would need to grow---the opposite of the missing type. The
spacing/packing lemma, the chain-class volume law, and the dilate theorem
were each examined as potential margin carriers and each resolves the
same way: a $k$-growing \emph{count} whose margin over its own trivial
baseline is $k$-free, a degrading order-factor slack, or a conditional
consequence of the hypothesis stack it was meant to prove.

The empirical half of the reading (``no proved margin-type statement
exists in the inventory'') is thereby settled. The structural half (``no
classification toolkit \emph{can} produce one'') is a hypothesis about
tool reach; it is what Instances 2--4 test, in three independent
settings, and what the falsity archive of \S\ref{sec:witness} completes
on the metric side.

\subsection{Instance 2: the transfer-matrix route, closed by theorem}
\label{sec:tm-transfer}

The one margin-shaped candidate instrument in the entire program was a
transfer-matrix (automaton) model of the solution cascade: transfer
operators naturally emit spectral quantities uniform in the scale
parameter, and a uniform spectral gap would be exactly the missing type.
The session that built it fixed its success criterion before computing:
the whole-cascade survival fraction must be bounded by
$\rho^{\mathrm{depth}}$ times the initial mass, with spectral radius
$\rho\le 1-\epsilon$ and $\epsilon$ independent of scale.

\begin{lemma}[Monomial dichotomy; Lemma M of the program records]
\label{lem:monomial}
Every constrained transfer matrix of the lossless cascade is a partial
permutation matrix. Its spectrum is $\{0\}$ together with roots of unity;
its spectral radius is exactly $0$ or $1$, never interior, at any scale,
in any state formulation. The affine pair maps compose to the identity
around any closed fiber loop (the dilation ratios telescope), so the
cyclic transfer product is diagonal, and its trace is the cascade's own
survivor count.
\end{lemma}

The proof is short once the state space is fixed: the cascade is
non-consuming (pure intersection semantics), the drift maps are
bijections, and a bijection constrained to an invariant subset is a
partial permutation. The machine record is not short: $1500/1500$ random
constrained systems; $12{,}000$ exact trace fingerprints over $300$
systems; $66$ homogenized pair systems on the committed families, with
spectral radius $1$ at exactly the six reflection-symmetric families
(their deletion sets close into $2$-cycles and the two-fiber survival
fraction is exactly $1.0000$ at all depths) and $0$ at the other sixty
(total extinction within three iterates, the intermediates set by lattice
runs of solution points, not by geometric decay). The criterion is
unattainable by theorem, not merely unattained.

Two further doors were checked. Coarse-grained or weighted automata can
carry interior spectral radii numerically, but the radius dominates every
cycle's geometric mean, so a uniform gap would force uniform decay of the
census weights---the margin as an \emph{input}, not an output---and the
converse fails (a $2\times2$ witness with all cycle means $\le0.9$ and
radius $1.122$: branching resurrects). Averaged or randomized operators
collapse to the independent-extension baseline. And the reduction closes
the loop: the whole-cascade survival fraction equals the product of
consecutive conditional ratios, each bounded by the pair-volume and
solution-set thinness hypotheses the route was commissioned to prove. The
instrument is the census in different notation; it consumes the
hypotheses it was built to generate. Losslessness forbids generated
decay---bijections preserve measure, deletion filters without
dissipating---and the reformulation is bounded by its own fidelity.

\subsection{Instance 3: character rigidity and saturated envelopes}
\label{sec:tm-char}

The harmonic-analytic route phrased the cascade obstruction in terms of
\emph{characters}: the four singleton coordinates and six pairwise
differences of the placement vector form a closed system of ten
functionals (Lemma C10 of the program records), closed under the
renormalization action up to affine shifts, whose aggregate
marginal-multiset is an orbit invariant. Any class-level census law must
live on this system; the route's natural strengthenings are laws about
its per-character profiles.

They are false, exactly. The singleton-profile orbit invariance (Lemma
V$'$) fails at both fully-enumerated primes: per-key counts match across
orbit members only after a size permutation, and the profiles differ. The
pairwise-profile invariance (V$''$) fails with $0/32$ sigma-matched
tuples. The mechanism is \emph{anchor mixing}: the renormalization
bijection couples every moving slot to the anchor coordinate; slot
singletons map to anchor pairs and pairwise differences to non-anchor
singletons, so no per-character statistic can be invariant---only the
aggregate is.

Worse, the per-character envelopes are \emph{saturated} exactly where the
margin would be needed. On the committed top orbits the pinning envelope
equals the modulus at both $k=2$ primes and $27/29$ at $k=3$ (fullness
fractions $1.000$, $1.000$, $0.931$): the marginal boxes are full, the
$\gamma$-envelope sits at or above the threshold value $4$ at $k=2$, and
the only invariant sparsity is \emph{joint} (median joint densities in
the marginal box: $1.5\%$, $7.4\%$, $6.9\%$). A bound phrased on
individual characters is tight---zero slack---precisely at the
configurations the obligations target.

Finally, the trivial bound sits exactly at the cascade threshold. The
baseline $|Sol|\le p^{m-1}$ realizes the constant
$c_1=(T/(T-6)+o(1))^{m-1}$---numerically $256$, $243$, $244$ at $T=8,9,10$
---which is exactly the value at which the depth denominator of the
conditional scaling theorem vanishes. The hypothesis H1g's content is to
beat the trivial bound by a factor growing in $k$; the measured factors
($95\times$, $1{,}207\times$, $3{,}803\times$) grow, but the proved ones
do not exist, and the sharp-constant question and the margin question
are one criterion stated twice. This is the type mismatch in its purest
form: the character system is exactly rich enough to locate the
obstruction and exactly too poor to bound it.

\subsection{Instance 4: the Lonely Runner zonotope}\label{sec:tm-zono}

The fourth instance is metric, and its falsity archive is the paper's
sharpest. The quotient-torus model of \S\ref{sec:prelim} gives two
statements: the $\delta$-form (the center class is within
$\mathrm{thr}_n$ of the lattice), which is equivalent to $\lrc_n$, and
the $\rho$-form (every class is), which is strictly stronger and would
open the geometry-of-numbers toolkit---transference bounds, covering
radius inequalities, dilution arguments. The natural lemma a
covering-radius toolkit wants first is exactly $\rho\le\mathrm{thr}_n$.

It is false, with exact witnesses, from $n=5$. At the extremal harmonic
instance $v=(1,2,3,4,5)$, the rational point
$c=(\tfrac1{32},\tfrac{15}{32},\tfrac3{32},\tfrac78)$ has
$D(c)=97/288>\tfrac13=\mathrm{thr}_5$: the deepest hole is not the
center, and the covering-radius statement fails at the very instance
where the $\delta$-form is tight. At $n=6$ the harmonic witness
$c=(0,\tfrac5{16},\tfrac{15}{16},\tfrac34,\tfrac12)$ gives
$D=29/80>\tfrac5{14}=\mathrm{thr}_6$, and the non-extremal instance
$v=(1,2,3,4,5,7)$ carries $c=(\tfrac1{16},\tfrac9{16},0,\tfrac38,0)$
with $D=4/11>\mathrm{thr}_6$: the falsity is not an extremal-curiosity
confined to the tight family. The rung-2 instance
$v=(1,2,3,4,5,12)$ has a witness at exactly $5/14=\mathrm{thr}_6$. All
witnesses are non-torsion points---the deepest holes are not the center
classes---and all are single exact evaluations of \eqref{eq:runnerform},
checkable by hand.

For $n\le4$ the two statements coincide on every instance tested: at
$n=3$ the deepest-hole property (that the maximum of $D$ is attained at
the center class) holds exactly when $3\mid m$ in the family
$(1,2,m)$ for $m\le24$, and at $n=4$ it holds when
$4\mid m$ for $m\le16$ (away-certified at $m=4,8,12$;
branch-and-bound certified at $m=16$, \S\ref{sec:bb}), with exact
witnesses against every other $m$.
The coincidence of the $\delta$- and $\rho$-forms in the small cases is
precisely why the stronger form was tempting; the $n=5$ witness is
precisely why it is wrong.

The companion hope---that the zonotope's Ehrhart theory could control
the covering radius through the positivity structure of its
$h^*$-polynomial---fails independently, and its failure mode is the
cleanest illustration of type mismatch in the whole program. The
$h^*$-polynomials are real-rooted, log-concave, and Newton, uniformly,
in every instance tested through $n=10$ in both core families
(Appendix~\ref{sec:hstar}). But the log-concavity \emph{strictness}
decays like $\Theta(1/n)$ ($n\cdot\text{strictness}$ drifting from
$\approx37$ to $\approx30$ over $n=3,\dots,10$), so there is no growing
quantitative margin to spend; and across the instance sweeps the
strictness has no measurable correlation with the covering margins
(Spearman $+0.32$ at $n=3$, $-0.31$ at $n=4$; the sign flips). A
coefficient inequality of a counting polynomial at growing scale $t$ is
the wrong type to pin a fixed-scale metric hole depth; the computation
says so twice, once by decay and once by correlation, and the
covering-radius statement it was meant to feed is false at the extremal
instances anyway. There is no mechanism, and none is needed: the bridge
is dead on both ends.

\subsection{The pattern}

Table~\ref{tab:instances} assembles the four instances. Each row names
the reformulation, the toolkit strengthening it wanted, the status of
that strengthening, and the witness. The rows are independent; the last
column is always finite, exact, and small. We emphasize what the table
does \emph{not} contain: any statement falsifying $\lrc$ itself, or any
statement weakening our certainty that the surviving equivalent forms
(the $\delta$-form, the pair-sum form) are exactly as hard as the
original.

\begin{table}[htbp]
\centering
\small
\caption{The four instances of type mismatch. ``Margin type'' is the
statement class the toolkit needed; ``status'' is its verified fate.}
\label{tab:instances}
\begin{tabularx}{\textwidth}{@{} l >{\raggedright\arraybackslash}X
>{\raggedright\arraybackslash}X l @{}}
\toprule
Reformulation & Toolkit wanted & Status & Witness \\
\midrule
Pair-sum covering & growth-in-$k$ margin & no proved instance in full
inventory; $k$-cancellation is a $k$-invariance & audit record
($34$ statements) \\
Transfer cascade & spectral gap $1-\epsilon$ & false by Lemma~M;
survival $=1$ at reflection families & six families, exact \\
Character rigidity & per-character profile laws & V$'$/V$''$ false;
envelopes saturated; trivial bound at threshold & exhaustive orbits;
$0/32$ \\
LR-zonotope & $\rho\le\mathrm{thr}_n$ (covering) & false from $n=5$;
$h^*$ strictness decays, no correlation & $97/288$, $29/80$, $4/11$
(exact) \\
\bottomrule
\end{tabularx}
\end{table}

% ======================================================================
% Section 4 — The positive core
% ======================================================================

\section{The positive core: classification at the rigid primes and at
\boldmath$p^2$/$pq$}\label{sec:positive}

The program's attacks were launched from a positive base, and that base
survives the negative results intact. This section states it with
precision; the full proofs, the verification records, and the
measurement methodology are the subject of the companion monograph
\cite{prog-artifacts}, and nothing here improves or weakens them. What
matters for this paper's thesis is that the positive core is exactly as
bounded as it is: everything it proves is a closure, classification, or
identity, and its one conditional theorem wears its hypotheses on its
sleeve.

\subsection{The sum-only equality theorem}

The analytic folklore (recorded in a comment on Tao's weblog
\cite{Tao-blog}) is that the optimal time is rational with denominator
the sum or difference of two speeds. The sharpening the program proved
is that \emph{sums alone suffice}.

\begin{theorem}[Sum-only equality; Theorem 1.1 of the companion
monograph]\label{thm:equality}
Let $V=(v_1,\dots,v_n)$ have distinct positive entries with
$\gcd(V)=1$ and $0<M(V)<\tfrac12$. For each unordered pair $\{u,w\}$
let $\tau(u,w)$ denote the optimal loneliness of the
$(n-1)$-runner instance obtained by restricting to the pair-sum lattice
of denominator $N=v_u+v_w$. Then
\[
M(V)\;=\;\max_{\{u,w\}}\tau(u,w),
\]
and consequently $\lrc_n$ holds for all $V$ if and only if the finite
statement $(\taust_n)$ does: for every $V$, some pair-sum lattice
carries an optimal value $\ge 1/(n+1)$. The finitization is exact; no
information is lost, and nothing about the classical difficulty is
claimed to become easier.
\end{theorem}

The two ingredients are that the binding runners at a global optimum can
be chosen two-sided, forcing the optimal time onto a \emph{sum} lattice
(difference denominators never carry the optimum), and a mirror lemma:
on that lattice the two binding runners are exact reflections of each
other at every lattice time, so the $n$-runner minimum there equals an
$(n-1)$-runner minimum. The equivalence of Corollary form
($\lrc_n\iff\taust_n$) is the lossless gateway through which every
instance of \S\ref{sec:tm-covering} entered.

\subsection{The covering framework and the rigid zone}

Under the equality theorem, a pair \emph{certifies} an instance unless
its $n-1$ bad sets $B_w=\{k\in\Z_N:(n{+}1)\wn{wk}<N\}$ cover $\Z_N$.
The framework that follows is finite and exact, and its proved pieces
include a size law for the bad sets, a polygonal witness lemma, moat
counting, a coincidence-counting lemma, a lift law, and a class-group
reduction at prime moduli (Lemma 2.2--2.7 of the companion
monograph). Five empirically distinct obstruction classes from the
program's early experiments collapse, in this language, into additive
orders of one covering phenomenon.

\begin{theorem}[Rigid-zone classifications; Theorems 3.1--3.5 of the
companion monograph]\label{thm:rigid}
At rung $n=4$ the unit-capable moduli---those $N$ for which $n-1$ unit
bad sets can cover $\Z_N$---are $7,11,13$; at rung $n=5$ they are
$7,13,17,19,37$ within the verified range (primes $\le150$, composites
$\le111$). At every one of these primes the covering configurations are
classified completely by a short explicit list of translate classes in
the cyclic group: antipodal classes and domino tilings at $7$ and $11$;
perfect matchings at $13$; the classes $\{0,1,2,7\}$, $\{0,2,4,6\}$ at
$17$ and $19$ (together with one further class at $19$); and the single
class $\{0,2,12,15\}$ at $37$.
\end{theorem}

This is the program's central structural finding: \emph{at rigid primes,
the covering obstruction is governed by a small translate-class
structure}, and the classification is what closes $75.6\%$ of the
$3{,}712{,}576$-instance corpus outright. The residual obstruction is
localized (mixed order signatures at composite moduli), and the composite
side carries the structure that the next subsection records.

\subsection{The composite side: group ring, tensor, and the fiber
cascade}

Beyond the unit capable moduli, at $N=p^2$ and $N=pq$, the program built
the composite-side structure that the scaling question lives on. The
group-ring translation expresses the covering count as a cyclotomic
coefficient problem; the $pq$ side decomposes as a tensor problem whose
Lam--Leung-type structure \cite{LamLeung} fixes the admissible
fiber-count signatures; and the foot-count input to the cascade was
discharged by an exact identity verified on $10{,}410$ triples with the
fixed threshold at $p=31$. The higher rungs carry a general-$T$ fiber
machinery, an exact dilate-intersection lemma, a reflection-fiber
invariance, a chain-class dilate theorem, and a renormalization lemma;
the Mirsky--Newman and de Bruijn--Schoenberg covering-system identities
\cite{MirskyNewman,deBruijnSchoenberg} play the classical supporting
roles.

Everything above is proved and machine-verified. What is \emph{not}
proved is the scaling closure itself, and the program consolidated it
into a single conditional statement with a disclosed ledger.

\begin{theorem}[Conditional scaling closure; Theorem 5.1 of the
companion monograph]\label{thm:SC}
Assume the distinct-family volume law H1g and the dilate-sparsity pair
H2m/H2r. Then the fiber cascade closes the covering obstruction at the
frontier cells $(1K,T{-}3U)$ of every rung $T\ge8$ at $N=p^2$, within a
depth that is $T$-flat in practice and never exceeds $T+3$ in the worst
case. The reduction chain and every other input are proved or
machine-validated; each constant's provenance---derived, measured, or
fitted, with the data it is fit to---is disclosed in the six-entry
hypothesis ledger, and the two obligations are reduced to named objects:
the strand census (for H1g) and the sharp pair-line constants (for
H2m).
\end{theorem}

We do not restate the ledger here. Its relevance to this paper is its
shape: the proved side of the ledger is all closures and identities; the
open side is exactly the two growth-in-$k$ margins whose type-mismatch
career occupies \S\ref{sec:tm}. Theorem~\ref{thm:SC} is what the positive
core amounts to when stated at its honest strength, and it is why the
program's claim to a positive result is deliberately the bounded one:
\emph{the obstruction is completely classified where classification
suffices, and the residual is a named margin of a type no turn of the
program ever proved}.

% ======================================================================
% Section 5 — The witness archive
% ======================================================================

\section{The witness archive}\label{sec:witness}

The archive has three shelves: exact refutations of the covering-radius
form (each a single rational point, checkable by one evaluation of
\eqref{eq:runnerform}), the deepest-hole certification ladder for the
family $(1,\dots,n{-}1,m)$, and the branch-and-bound certificate method
introduced in this session, together with its negative-control
validation. Every entry is re-asserted in the provenance run of
Appendix~\ref{sec:prov}.

\subsection{Exact refutations of the covering-radius form}

Table~\ref{tab:refute} lists the witnesses. Each witness $c$ is a
rational point of the quotient torus in $c$-coordinates; the claimed
value is $D(c)$ from \eqref{eq:runnerform}, computed exactly and
independently by the linear-form/lattice-search implementation of the
program's records. The harmonic instances are extremal for the
$\delta$-form ($\delta=\mathrm{thr}_n$ exactly), so these witnesses kill
the $\rho$-form at the very instances where the surviving equivalent form
is tight; the $n=6$, $m=7$ witness kills it at a generic, non-extremal
instance; and the rung-2 instance at $m=12$ sits exactly at the
threshold. The witnesses are non-torsion in every case: the deepest
holes are not the center classes.

\begin{table}[htbp]
\centering
\footnotesize
\caption{Exact witnesses against the covering-radius form
$\rho(V)\le\mathrm{thr}_n$. All values are exact rational lower bounds
for $\rho$; enclosures from the grid searches are recorded in the
program's data. Witness points are in $c$-coordinates
(\S\ref{sec:prelim}).}
\label{tab:refute}
\begin{tabularx}{\textwidth}{@{} l l l >{\raggedright\arraybackslash}X l @{}}
\toprule
$V$ & $\delta=\tfrac12-\lambda$ & $\mathrm{thr}_n$ & witness $c$ &
$D(c)$ \\
\midrule
$(1,2,3,4,5)$ & $1/3$ & $1/3$ &
$(\tfrac1{32},\tfrac{15}{32},\tfrac3{32},\tfrac78)$ &
$97/288{\approx}.3368$ \\
$(1,2,3,4,5,6)$ & $5/14$ & $5/14$ &
$(0,\tfrac5{16},\tfrac{15}{16},\tfrac34,\tfrac12)$ & $29/80=.3625$ \\
$(1,2,3,4,5,7)$ & $1/3$ & $5/14$ &
$(\tfrac1{16},\tfrac9{16},0,\tfrac38,0)$ & $4/11{\approx}.3636$ \\
$(1,2,3,4,5,12)$ & $9/26$ & $5/14$ & non-torsion point &
$5/14=\mathrm{thr}_6$ \\
\bottomrule
\end{tabularx}
\end{table}

Three further exact refutation families fill out the record: the $n=5$
family $(1,2,3,4,m)$ fails deepest-hole for every non-multiple of $5$
up to $m=30$ (twenty exact witnesses, the sharpest being $6/19>5/16$ at
$m=15$); the $n=6$ family fails for every tested $m\in\{6,\dots,12,18,
24\}$, including all multiples of $6$ through $4n$ (extinction of the
modular deepest-hole law at $n=6$; the narrowest margin is $7/5200$ at
$m=24$, witness value $71/208$); and at $n=6$, $m=9$ the witness sits at
$\rho\ge\mathrm{thr}_6$ exactly. A useful way to read the table: the
$\rho$-form was already falsified three times over before this session,
and the session's contribution below is on the \emph{positive} side of
the ladder---certifying where the deepest hole genuinely is the center.

\subsection{The deepest-hole certification ladder}

Table~\ref{tab:ladder} records the corrected modular law with its
certification levels. The law as originally stated (``deepest hole at
the center iff $n\mid m$'') is a small-$n$ phenomenon, genuine and exact
at $n=3$, certified at $n=4$, broken at $n=5$ by the harmonic instance
itself, and extinct at $n=6$; the record below is the version that
survived the audit, and the two rows improved by this session's
branch-and-bound method are marked.

\begin{table}[htbp]
\centering
\small
\caption{The deepest-hole ladder for $(1,\dots,n{-}1,m)$: when is the
maximum of $D$ attained at the center class ($\rho=\delta$)?
``Exact'' = complete exact proof; ``certified'' = complete proof with an
independent outer loop; ``consistent'' = search-level only.}
\label{tab:ladder}
\begin{tabularx}{\textwidth}{@{} c c >{\raggedright\arraybackslash}X @{}}
\toprule
$n$ & tested $m$ & status of ``deepest hole at center iff $n\mid m$''
\\
\midrule
$3$ & $3..24$ & exact: holds iff $3\mid m$; exact witnesses against all
others \\
$4$ & $4..16$ & certified: holds iff $4\mid m$ ($m=4,8,12$
away-certified; $m=16$ certified by branch and bound, \emph{this
session}); exact witnesses against all non-multiples \\
$5$ & $5..30$ & law false: $m=5$ (harmonic) and $m=15$ fail exactly;
$m=20,25,30$ certified by branch and bound (\emph{this session});
$m=10$ open after one serious attempt (\S\ref{sec:bb}); all $20$
non-multiples fail exactly \\
$6$ & $6..12,18,24$ & extinct: every tested $m$ fails exactly, including
every multiple of $6$ through $4n$ \\
\bottomrule
\end{tabularx}
\end{table}

The open residue is the $n=5$ multiples that survived the audit:
$m=10,20,25,30$ carry the hold signature (the torsion center attains
$\delta$; the $K=32$ grid maximum equals $\delta$ exactly; a flat ridge
passes through the center; a $288$-start Nelder--Mead search with exact
re-evaluation finds nothing above $\delta$), but at $d=4$ a global exact
certificate was beyond the program's earlier arrangement machinery.
This session's branch-and-bound method (\S\ref{sec:bb}) upgrades three
of the four. At $m=25$ and $m=30$ the trees empty in under five minutes
each ($10{,}414$ and $10{,}646$ certified leaves; $\rho=4/13$ and
$\rho=19/62$ exactly), and at $m=20$ the tree empties in sixteen
minutes ($37{,}018$ leaves; $\rho=13/42$). At $m=10$ the single
serious attempt ($25$ minutes, $143{,}228$ box evaluations, $61{,}038$
exact leaf certificates) did not terminate: the unresolved set has
total volume below $6\times10^{-7}$, the exact depths at its sampled
centers are at most $1135/3584<7/22$, and the largest exact $U$-bounds
among the sampled survivors are all exactly
$571/1792=7/22+9/19712$---a quantized minimax loss, not an excess. The
status of $m=10$ is therefore: \emph{consistent, with the two-sided
evidence substantially strengthened, and exact equality open at $d=4$}.
We record it as the archive's one deliberately unfinished row; the odd
residue ($k\in\{1,3\}$ failing, $k=2$ open, $k\in\{4,5,6\}$
certified) is noted as an open pattern rather than forced into a law.

\subsection{The branch-and-bound certificate and its negative control}
\label{sec:bb}

This session introduced a second, independent certification method,
built on the runner formulation \eqref{eq:runnerform} rather than on
the arrangement enumeration of the program's records. For a dyadic
$c$-box $C$ define
\[
U(C)\;=\;\min_{\sigma\in[0,1)}\;\max_j\;\sup_{c\in C}\wnc{c_j-v_j\sigma},
\]
computed exactly: each inner supremum is piecewise linear in $\sigma$
with a rational kink set, so the minimum of their maximum is attained at
a rational kink or crossing. By the minimax inequality,
$\max_C D\le U(C)$, so a box with exact $U(C)\le r$ is \emph{certified}
to contain no point of depth above $r$; boxes that fail are split
dyadically. If the tree empties, $\rho\le r$ is exactly certified by a
finite, auditable leaf list. Two soundness refinements matter in
practice: the $\sigma$-search is restricted to
$[0,r]\cup[1-r,1)$, which is lossless because any certifying $\sigma$
has $\wnc\sigma\le r$; and equality $U(C)=r$ prunes, which the tight
boxes need.

The method carries its own negative control, and the control behaves
exactly as it must. Run on the harmonic instance $v=(1,2,3,4,5)$ at
$r=\delta=\tfrac13$---where the truth is $\rho\ge97/288>\tfrac13$---the
branch and bound \emph{refuses to certify}: its survivor set persists
around the true deep holes, and the exact depths at the survivors'
centers reach $607/1792\approx0.3388>\tfrac13$, at
$c=(\tfrac{11}{256},\tfrac{163}{256},\tfrac{15}{256},\tfrac{79}{256})$
---a marginally deeper witness than the recorded $97/288$, found
incidentally by the box centers and re-asserted exactly in the follow-up
run, which improves the harmonic lower bound to $\rho\ge607/1792$. The machinery thus demonstrates both halves of its own
soundness profile on one instance: it certifies only what is true, and
it localizes the false. The same control distinguishes the two
$n=5$ worlds sharply: for the harmonic instance the surviving excess
stays above the threshold by a fixed margin, while for the consistent
rung-2 instances the surviving excess is pure minimax loss and decays
with the box size.

\subsection{Summary of the archive}

The archive's content, stated as a single paragraph: the
covering-radius strengthening of the Lonely Runner statement is false
from $n=5$, with three exact witnesses at extremal and non-extremal
instances, an improved harmonic lower bound $\rho\ge607/1792$, and an
extinction record at $n=6$; the equivalent
center-depth form survives untouched (it is $\lrc$); the deepest-hole
modular law is exact at $n=3$, certified at $n=4$ (both sides, all
tested $m$, including $m=16$ by the new method), false with one open
residue at $n=5$ (the $k=1,3$ multiples fail exactly; $k=4,5,6$ are now
certified; $k=2$ is open after a serious attempt), and extinct at $n=6$; and every number in the paragraph is
re-asserted by a persisted, independent, exact computation. The archive
is finite, small, and open to inspection; that is the point.

% ======================================================================
% Section 6 — Conclusion
% ======================================================================

\section{Conclusion}\label{sec:concl}

\subsection{What is established}

Four facts, each carrying its own verification record. First, the
covering-radius form of the Lonely Runner statement---the form that
would have opened the geometry-of-numbers toolkit---is strictly stronger
than the conjecture and false from $n=5$, with exact rational witnesses
at extremal and non-extremal instances and an extinction record at
$n=6$. Second, the transfer-matrix route into the strand census is
closed by theorem: the lossless cascade's constrained dynamics has no
interior rate at any scale, and the survival fraction is exactly $1$ at
the reflection-symmetric families. Third, the character-rigidity route's
natural strengthenings are false with exhaustive exact witnesses, and
its marginal envelopes are saturated exactly where the missing margin
would have to live. Fourth, on the classification side, the full proof
inventory of the program's covering reformulation contains no proved
statement of the margin type---verified statement-by-statement---and the
one instrument shaped to produce that type is the one the second fact
closes.

Against these stands the bounded positive core: the sum-only equality
theorem and its lossless finite reformulation; the rigid-zone
translate-class classifications at the primes $7$, $11$, $13$, $17$,
$19$, $37$; the composite-side structure at $p^2$ and $pq$; the
conditional scaling theorem with its disclosed hypothesis ledger; and,
from this session, an independently certified deepest-hole archive
(equality $\rho=\delta$ at $n=4$, $m=16$ and at $n=5$,
$m=20,25,30$, with the improved harmonic lower bound
$\rho\ge607/1792$ and the $m=10$ residual characterized as pure
minimax loss). And
beneath both, the mechanism that ties the paper together: losslessness
preserves hardness exactly, toolkit strengthenings are false where they
matter, and the instruments a faithful reformulation supports compose
only into closures. The four instances are four independent
confirmations of the same three-part statement.

\subsection{Why the reformulations cannot cross}

It is worth saying plainly what the mechanism predicts, since the
prediction is falsifiable and was falsifiable at each step. Any future
attack on $\lrc$ through a lossless reformulation will face the same
fork. On one branch, the attack respects the equivalence and uses only
structure invariant under the rewriting; then it is an attack on
$\lrc$ itself, with the reformulation as notation, and its difficulty
is the conjecture's own. On the other branch, the attack strengthens
the reformulation---asks for a margin, a gap, a rate, a monotone law---
and the strengthening is a strictly stronger statement than $\lrc$;
at the extremal instances it is false, and the falsity is exhibited by
small exact witnesses. The fork is not a metaphor: in each of the four
instances the program went down one branch, hit exactly the predicted
wall, and recorded the wall's coordinates.

The prediction for the zonotope's residual positive question is the same
prediction with the sign of interest reversed. Real-rootedness of the
$h^*$-polynomials of the LR-zonotopes---verified through $n=10$ in both
core families and in every sweep instance tested---may well be a true
theorem of the arithmetic-matroid type, and if it is, it belongs to the
arithmetic-matroid literature, decoupled from the Lonely Runner problem
(Appendix~\ref{sec:hstar}). What it cannot be, given the decay and
non-correlation records of Appendix~\ref{sec:hstar} and the falsity of
the covering-radius statement it was meant to feed, is a bridge to
$\lrc$. That was always the fork's second branch.

\subsection{What would be needed}

The program's records end with a scope statement rather than a
speculation, and we keep to it. A genuine advance on $\lrc$ through any
of the channels this paper surveys would require an instrument whose
output type is a uniform margin over a trivial baseline, independent of
the scale parameter, at the extremal instances where every natural
candidate saturates or vanishes. No such instrument appears in the
covering toolkit, the spectral toolkit, the character toolkit, or the
Ehrhart toolkit, and the four closures explain structurally why the
obvious candidates from each cannot be one. Whether such an instrument
exists anywhere---or whether the conjecture's resistance is exactly this
fork, everywhere---is a question this paper does not answer and does
not need to answer to justify its archive.

\subsection{Scope, stated once more}

This paper proves no new cases of the Lonely Runner Conjecture. It does
not reduce, weaken, or reformulate-away the conjecture; the surviving
equivalent forms are as hard as the original, and the paper's
equivalence results make them no easier. The falsified strengthenings
were, by construction, stronger than the conjecture, so their falsity
leaves the conjecture exactly as it was: open, tight at the harmonic
instances, and now surrounded on four sides by precisely mapped walls.
The walls are the contribution. They are finite, exact, and checkable,
and they are documented so that the next program does not spend its
years rediscovering them.

% ======================================================================
% Appendix A — h* observations
% ======================================================================

\appendix
\section{The \boldmath$h^*$-observations through \boldmath$n=10$}
\label{sec:hstar}

This appendix records the Ehrhart observations on the LR-zonotopes,
extended by this session to $n=9$ and $n=10$ in both core families. The
extension was executed as a one-day addendum with a narrow purpose:
establish whether the real-rootedness and log-concavity observed at
$n\le8$ continue, and whether the strictness decay and the
non-correlation with covering margins persist. Both answers are yes. The
appendix closes with the statement we regard as the residual question of
independent interest.

\subsection{The table}

Table~\ref{tab:hstar} lists, for each family and $n$: the
$h^*$-polynomial, real-rootedness (verified at $45$ decimal places on
the exact integer coefficients, with a companion-matrix check agreeing),
log-concavity and Newton's inequalities (exact rational arithmetic), and
the minimum ratio $h_k^2/(h_{k-1}h_{k+1})$, which is the
\emph{strictness} of log-concavity. Every row's Ehrhart data was
verified three ways (Eulerian expansion, standard inversion, and an
independent depth-first coset count at $t=1,2$), the volume identity
$e_{n-1}=\sum v_j$ was asserted exactly, and the loneliness identities
$\lambda=1/(n{+}1)$ (harmonic) and $\lambda=2/(2n{+}1)$ (rung-2) were
re-asserted exactly. The $n\le8$ rows reproduce the program's earlier
records digit for digit; the $n=9,10$ rows are new.

\begin{table}[htbp]
\centering
\footnotesize
\caption{$h^*$-polynomials of the LR-zonotopes, $n=3..10$, both core
families. ``RR'' = real-rooted (exact-integer verification at 45 dps);
``LC/Nw'' = log-concave and Newton (exact rational arithmetic);
$S$ = $\min_k h_k^2/(h_{k-1}h_{k+1})$; $nS$ = $n\cdot S$. The $n\le8$
rows reproduce the program's earlier records digit for digit; the
$n=9,10$ rows are new (this session).}
\label{tab:hstar}
\begin{tabularx}{\textwidth}{@{} l c >{\raggedright\arraybackslash}X c c r r @{}}
\toprule
family & $n$ & $h^*(t)$ & RR & LC/Nw & $S$ & $nS$ \\
\midrule
harm. & 3 & $1+7t+4t^2$ & y & y & $12.250$ & $36.8$ \\
rung2 & 3 & $1+11t+6t^2$ & y & y & $20.167$ & $60.5$ \\
harm. & 4 & $1+18t+35t^2+6t^3$ & y & y & $9.257$ & $37.0$ \\
rung2 & 4 & $1+22t+51t^2+10t^3$ & y & y & $9.490$ & $38.0$ \\
harm. & 5 & $1+37t+179t^2+133t^3+10t^4$ & y & y & $6.511$ & $32.6$ \\
rung2 & 5 & $1+45t+239t^2+181t^3+14t^4$ & y & y & $7.013$ & $35.1$ \\
harm. & 6 & $1+78t+744t^2+1286t^3+399t^4+12t^5$ & y & y & $5.518$
& $33.1$ \\
rung2 & 6 & $1+86t+920t^2+1682t^3+535t^4+16t^5$ & y & y & $5.748$
& $34.5$ \\
harm. & 7 & $1+147t+2522t^2+8948t^3+7323t^4$\newline
$+1201t^5+18t^6$ & y & y & $4.335$ & $30.3$ \\
rung2 & 7 & $1+161t+3024t^2+11144t^3+9295t^4$\newline
$+1551t^5+24t^6$ & y & y & $4.418$ & $30.9$ \\
harm. & 8 & $1+290t+8317t^2+51458t^3+82947t^4$\newline
$+35270t^5+3135t^6+22t^7$ & y & y & $3.791$ & $30.3$ \\
rung2 & 8 & $1+298t+9277t^2+60986t^3+102275t^4$\newline
$+44798t^5+4095t^6+30t^7$ & y & y & $3.829$ & $30.6$ \\
harm. & 9 & $1+559t+25519t^2+258901t^3+734059t^4$\newline
$+631993t^5+155305t^6+8035t^7+28t^8$ & y & y & $3.293$ & $29.6$ \\
rung2 & 9 & $1+585t+28885t^2+305035t^3+880149t^4$\newline
$+764143t^5+188659t^6+9789t^7+34t^8$ & y & y & $3.323$ & $29.9$ \\
harm. & 10 & $1+1098t+77550t^2+1213626t^3+5529648t^4$\newline
$+8309694t^5+4180482t^6+626910t^7+19359t^8+32t^9$ & y & y & $2.987$
& $29.9$ \\
rung2 & 10 & $1+1114t+83790t^2+1377794t^3+6457600t^4$\newline
$+9866694t^5+5017826t^6+758710t^7+23631t^8+40t^9$ & y & y & $3.004$
& $30.0$ \\
\bottomrule
\end{tabularx}
\end{table}

The sweep families behave identically: every instance of $(1,2,m)$,
$m\le14$, and $(1,2,3,m)$, $m\le12$, tested by the program was
real-rooted, log-concave, and Newton.

\subsection{Strictness decay and non-correlation}

Two quantitative facts were already in the program's records at $n\le8$
and persist through the new rows. First, the strictness $S$ decays
monotonically in $n$ in both families, with $n\cdot S$ drifting from
$\approx37$ at $n=3,4$ to $\approx29.6$--$30.0$ at $n=9,10$: the decay
is $\Theta(1/n)$ with a slowly drifting constant, so there is no growing
quantitative margin in the coefficient structure to spend. Second, the
strictness carries no correlation with the covering margins: across the
$n=3$ sweep (twelve instances, $(1,2,m)$, $m=3..14$) the Spearman
coefficient of strictness against the margin is $+0.32$; across the
$n=4$ sweep (eight instances, $(1,2,3,m)$, $m=5..12$) it is $-0.31$.
The sign flips between adjacent rungs; the harmonic family has margin
identically $0$ at every $n$ while its strictness varies by a factor of
four over the table. Both correlations were recomputed in the provenance
run from the persisted instance data and reproduce exactly.

The two facts close the bridge in the coefficient direction, and the
falsity of the covering-radius statement (\S\ref{sec:tm-zono}) closes it
in the metric direction: the statement that strict log-concavity would
have needed to feed is false at the extremal instances, and the
coefficient structure cannot see the metric hole depth that falsifies
it. There is no mechanism, and the arithmetic is clean enough that the
absence of the mechanism is itself a datum.

\subsection{The residual question}

What survives is a question of independent arithmetic interest, and we
state it decoupled from the Lonely Runner problem entirely. The
Ehrhart data of the LR-zonotopes is governed by the gcd structure
\eqref{eq:ehr} of the projected basis---the arithmetic matroid of the
configuration \cite{dAMoci}. In every instance tested (two core
families through $n=10$, sweeps through $m\le14$ at $n=3$ and $m\le12$
at $n=4$), the arithmetic $h^*$-polynomial is real-rooted. The natural
conjecture-shaped question is whether this holds for the projected
bases of the form $(1,\dots,n-1,m)$ or beyond, and the natural home for
it is the real-rootedness and Lorentzian literature for arithmetic
matroid Tutte specializations \cite{BrandenHuh,dAMoci,StanleyEC}. We
flag it as a question, record the data, and deliberately do not pursue
it further here: the extension to $n=9,10$ was executed to complete the
record, and the record's reading is that the question is real but
belongs to a different subject. It is not a bridge, and treating it as
one would repeat the pattern this paper documents.

% ======================================================================
% Appendix B — Computational provenance
% ======================================================================

\section{Computational provenance}\label{sec:prov}

The audit standard adopted for this paper is stated in one sentence:
\emph{every number cited in the text traces to a script that was
actually run and whose output was persisted.} It was applied
retroactively to all inherited numbers and prospectively to all new
ones. This appendix is the map.

\subsection{The ledger}

The provenance run (\texttt{scripts/paper\_provenance.py}; output
\texttt{out\_paper\_provenance.json} and the human-readable
\texttt{out\_paper\_provenance.txt}) reads the two persisted run records
of the zonotope program and re-asserts, by independent exact
computation, every value the paper cites from them: all witness depths
(evaluating the runner formulation \eqref{eq:runnerform} \emph{and} the
independent linear-form/lattice-search implementation, asserting exact
agreement of the two), all recorded $\rho$ values at their recorded
argmax points, all $\lambda$ and $\delta$ values, both Spearman
coefficients, both correlation instance arrays (strictness recomputed
from the Ehrhart coefficients and margins recomputed from exact
$\lambda$), and the strictness entries of the $n\le8$ table. The
complete ledger contains $225$ items. The number of failures is $0$.

The two-way witness verification deserves emphasis because it is the
strongest form of check the program has: every witness value in
Tables~\ref{tab:refute} and the fail-witness lists of
\S\ref{sec:witness} was recomputed by two implementations with different
mathematics (a one-dimensional minimization over the trajectory
parameter, and a pruned lattice search in the quotient norm's linear
forms), agreeing exactly. A one-line reproduction of the headline
witnesses, for the skeptical reader with the runner formulation at
hand: evaluate
$\min_\sigma\max(\wnc\sigma,\wnc{c_2-2\sigma},\wnc{c_3-3\sigma},
\wnc{c_4-4\sigma},\wnc{c_5-5\sigma})$ at
$c=(\tfrac1{32},\tfrac{15}{32},\tfrac3{32},\tfrac78)$; the minimum is
attained at $\sigma=\tfrac{31}{288}$ and the value is $97/288$.

\subsection{The new runs of this session}

Three scripts were written and run for this paper; their outputs are
persisted alongside the inherited records and copied into the paper's
artifact bundle.

\begin{enumerate}[leftmargin=1.6em,itemsep=2pt]
\item \texttt{hstar\_n9\_n10.py} $\to$ \texttt{out\_hstar\_n9\_n10.json}:
the $h^*$ extension of Appendix~\ref{sec:hstar} ($n=3..10$, both
families; Ehrhart asserted three ways; roots at $45$ decimal places on
exact integer coefficients; strictness, Newton, log-concavity exact;
$\lambda$ asserted against the classical values).
\item \texttt{certify\_rung2\_n5.py} $\to$
\texttt{out\_certify\_rung2.json}: the branch-and-bound certificate of
\S\ref{sec:bb}, including the cross-validation of the two $D$
implementations (forty random rational points per instance, exact
agreement), the controls ($n=3$ and $n=4$ core instances, all
certified), the negative control (harmonic $n=5$, correctly refusing),
and the ladder targets (\S\ref{sec:witness}).
\item \texttt{paper\_provenance.py} $\to$
\texttt{out\_paper\_provenance.\{json,txt\}}: the $225$-item ledger
above.
\item \texttt{certify\_followup.py} $\to$
\texttt{out\_certify\_followup.json}: the survivor enclosure analysis
of the open $m=10$ run and the recovered improved harmonic witness
($607/1792$ at $c=(\tfrac{11}{256},\tfrac{163}{256},
\tfrac{15}{256},\tfrac{79}{256})$).
\end{enumerate}

\subsection{Certificates and budgets}

The branch-and-bound certificates are finite objects: a certificate for
$\rho(V)\le r$ is the pruned dyadic tree, each leaf carrying an exact
rational $U$-value $\le r$ computed from the kink/crossover enumeration
of \S\ref{sec:bb}. Budgets and certificate sizes for the runs cited in
the text are listed in Table~\ref{tab:certs}; the JSON persists leaf
samples, survivor aggregates (count, bounding box, maximum excess over
$r$, exact depths at sampled survivor centers), and the exact
$U$-values, so every certificate can be re-audited without re-running
the search, and re-running the search reproduces it deterministically.

\begin{table}[htbp]
\centering
\small
\caption{Branch-and-bound runs cited in the text. ``leaves'' = certified
boxes (each with exact $U\le r$); ``surv.'' = unresolved boxes at cap;
wall-clock seconds in parentheses.}
\label{tab:certs}
\begin{tabular*}{\textwidth}{@{\extracolsep{\fill}} l l l l l}
\toprule
instance & $r$ & verdict & leaves & surv. \\
\midrule
$(1,2,3)$ & $1/4$ & certified & $14$ $(0.0)$ & $0$ \\
$(1,2,6)$ & $3/14$ & certified & $20$ $(0.0)$ & $0$ \\
$(1,2,3,4)$ & $3/10$ & certified & $438$ $(2)$ & $0$ \\
$(1,2,3,8)$ & $5/18$ & certified & $156$ $(1)$ & $0$ \\
$(1,2,3,12)$ & $7/26$ & certified & $168$ $(1)$ & $0$ \\
$(1,2,3,16)$ & $9/34$ & certified & $222$ $(2)$ & $0$ \\
$(1,2,3,4,5)$ & $1/3$ & \emph{refused} (control) & --- &
$45{,}355$ $(421)$ \\
$(1,2,3,4,10)$ & $7/22$ & open (one serious attempt) & $61{,}038$
$(1501)$ & $21{,}153$ \\
$(1,2,3,4,20)$ & $13/42$ & certified & $37{,}018$ $(981)$ & $0$ \\
$(1,2,3,4,25)$ & $4/13$ & certified & $10{,}414$ $(231)$ & $0$ \\
$(1,2,3,4,30)$ & $19/62$ & certified & $10{,}646$ $(274)$ & $0$ \\
\bottomrule
\end{tabular*}
\end{table}

\subsection{Reproduction}

All scripts run under CPython~3.12 with \texttt{fractions},
\texttt{numpy}, \texttt{scipy}, and \texttt{mpmath}, and take minutes
to hours on one core; the certificates of Table~\ref{tab:certs} take
seconds to minutes (the longest, $m=20$, took sixteen minutes), and the
refused control and the open $m=10$ attempt were run to their stated
budgets. The artifact bundle accompanying the paper contains the three
scripts, their outputs, the inherited zonotope records
(\texttt{out\_lrc\_zono\_final.json},
\texttt{out\_lrc\_zono\_rescan.json}, and the library
\texttt{lrc\_zono\_lib.py} they import), and the ledger. Nothing in the
paper relies on an unlogged computation, a hard-coded constant, or a
claim of the form ``verified elsewhere'': the audit that motivated this
standard found exactly those three failure modes in the program's own
earlier records, and the standard is the response.


% ======================================================================
% Bibliography (manual, ordered by first citation)
% ======================================================================

\begingroup\sloppy\raggedright
\begin{thebibliography}{16}

\bibitem{PS-survey}
G.~Perarnau and O.~Serra,
\emph{The Lonely Runner Conjecture turns 60},
arXiv:2409.20160 (survey).

\bibitem{Wills}
J.~M. Wills,
\emph{Zwei S\"atze \"uber inhomogene diophantische Approximation von
Irrationalzahlen},
Monatsh.\ Math.\ 71 (1967).

\bibitem{Rosenfeld8}
M.~Rosenfeld,
\emph{The lonely runner conjecture holds for eight runners},
arXiv:2509.14111 (2025).

\bibitem{Trakulthongchai}
T.~Trakulthongchai,
\emph{Nine and ten lonely runners},
arXiv:2511.22427 (2025).

\bibitem{Sungkawichai}
T.~Sungkawichai and T.~Trakulthongchai,
\emph{Eleven, twelve, and thirteen lonely runners},
arXiv:2604.23906 (2026).

\bibitem{Allikvere}
J.~Allikvere,
\emph{Fourteen and fifteen lonely runners},
arXiv:2609.02604 (2026).

\bibitem{MSS}
R.~D. Malikiosis, F.~Santos, and M.~Schymura,
\emph{Linearly-exponential checking is enough for the Lonely Runner
Conjecture and some of its variants},
arXiv:2411.06903 (2024).

\bibitem{Rifford}
L.~Rifford,
\emph{On the time for a runner to get lonely},
arXiv:2111.13688 (2021).

\bibitem{Tao-blog}
T.~Tao,
\emph{A remark on the lonely runner conjecture},
weblog post, 13 May 2015; the observation that the optimal time is
rational with denominator the sum or difference of two speeds appears
in the comments.
\\[2pt]
\url{https://terrytao.wordpress.com/2015/05/13/}\\
\url{a-remark-on-the-lonely-runner-conjecture}

\bibitem{prog-artifacts}
LRC reformulation program,
\emph{Companion monograph: The Lonely Runner Problem on Pair-Sum
Lattices---a lossless finitization and a conditional scaling-closure
theorem}, and the committed scripts, run logs, and experiment notes
referenced there; unpublished artifact bundle accompanying this paper
(2026).

\bibitem{BeckRobins}
M.~Beck and S.~Robins,
\emph{Computing the Continuous Discretely: Integer-Point Enumeration
in Polyhedra}, second ed.,
Undergraduate Texts in Mathematics, Springer, New York, 2015.

\bibitem{dAMoci}
L.~d'Adderio and L.~Moci,
\emph{Arithmetic matroids, the Tutte polynomial and toric
arrangements},
Adv.\ in Appl.\ Math. (2014).

\bibitem{StanleyEC}
R.~P. Stanley,
\emph{Enumerative Combinatorics, vol.~1}, second ed.,
Cambridge Studies in Advanced Mathematics 49, Cambridge Univ.\ Press,
2012.

\bibitem{Shephard}
G.~C. Shephard,
\emph{Combinatorial properties of associated zonotopes},
Canad.\ J.\ Math.\ 26 (1974), 302--321.

\bibitem{Frobenius}
G.~Frobenius,
\emph{\"Uber die Bernoullischen Zahlen und die Eulerschen Polynome},
Sitzungsber.\ K\"onigl.\ Preu\ss.\ Akad.\ Wiss.\ (1910).

\bibitem{LamLeung}
T.~Y.~Lam and K.~H.~Leung,
\emph{On the cyclotomic polynomial $\Phi_{pq}(X)$},
Amer.\ Math.\ Monthly \textbf{107} (2000), no.~7, 605--612.

\bibitem{MirskyNewman}
L.~Mirsky and D.~J.~Newman,
\emph{A problem in the theory of numbers},
Michigan Math.\ J.\ \textbf{9} (1962), 201--206; see also
H.~Davenport and R.~Rado, \emph{Covering systems of negative
congruences},
J.\ London Math.\ Soc.\ \textbf{38} (1963), 509--516.

\bibitem{deBruijnSchoenberg}
N.~G.~de~Bruijn and T.~J.~Schoenberg,
\emph{On the two vanishing identities of Euler},
Indag.\ Math.\ \textbf{23} (1961), 542--547.

\bibitem{BrandenHuh}
P.~Br\"anden and J.~Huh,
\emph{Lorentzian polynomials},
Ann.\ of Math.\ (2) \textbf{192} (2020), no.~3, 821--915.

\end{thebibliography}
\endgroup


\end{document}
